% Transcription — fonds Alexandre Grothendieck, shelfmark 13, batch 1 (pages 1-20).
% Established from the facsimile archives/batches/13/batch-01.pdf (a mirror of
% https://grothendieck.umontpellier.fr/13.pdf), per the skill
% transcribe-grothendieck. The batch file carries no cover sheet: PDF page N is
% archivists' page N. Checked against the pencilled numbers bottom-left — a "2"
% on the second leaf, then "12", "13", "15", "16", "17", "18", "19" on the
% corresponding leaves, confirm the alignment end to end.
%
% Pass: Opus 5 (claude-opus-5), 2026-09-19 — first pass, unchecked against the pages by a human.
%
% What the batch is. Two things, and a cover for a third. Page 1 is an
% otherwise blank leaf carrying, in his hand and underlined, « Groupe de Hodge »
% and « lettre à Mumford »: it is the wrapper for everything that follows, and
% it gets no \page. Pages 3 to 11 are nine leaves of a continuous French run, in
% his hand and unpaginated by him, on the Hodge group: the Deligne torus S, the
% homomorphisms j : S → G_R, polarisations, and the moduli space of such j. Page
% 12 is blank. Pages 13 and 14 are a letter in English to Mumford, dated on the
% leaf « Algiers Nov. 1, 1965 » — the letter of 1965 the inventory announces —
% and pages 15 to 19 are its long mathematical postscript, also in English, on
% the modular varieties coming from motivic Galois theory. Page 2 is blank apart
% from the show-through of page 3 and gets no \page. Page 20 is a second cover
% in his hand, « Descriptions transcendantes de motifs en car. 0 … », naming the
% run that begins in batch 2; it too gets no \page, and a \note records it. Its
% words are NOT made a section heading here: the run they name has no leaf in
% this batch, and batch 2 opens under that heading. Setting it in both files
% would print the same cover twice when the folder is read whole.
%
% The argument, page by page:
%
%    3    G algebraic over Q; S = ∏_{C/R} G_{m,C} with S(R) ≃ C*, and
%         U ⊂ S the norm-one subtorus. Rep(S) as real vector spaces with a
%         bigrading of V_C satisfying V^{p,q} = conj(V^{q,p}); λ acts on V^{p,q}
%         by λ^p λ̄^q. A functorial such bigrading on the forgetful composite
%         Rep(G) → Modf(Q) → Modf(R) → Modf(C) is the same as a homomorphism
%         j : S → G_R; a central i : G_m → G gives the total grading and the
%         compatibility V^n_C = ⊕_{p+q=n} V^{p,q}.
%    4    that compatibility is condition (a) j|G_{m,R} = i_R; Q(-1) as an
%         invertible object, i.e. ε : G → G_m; condition (b) εj = Norme, which
%         is what puts Q(1) in bidegree (1,1). Then the question: for V of
%         degree p, which forms φ : V × V → Q(-p) polarise the Hodge structure?
%         Answer (**): φ_R(x, Cy) positive definite.
%    5    C is √−1 ∈ S(R) ≃ C*, whence φ alternating in odd weight and
%         symmetric in even. G' = Ker ε and the recovery of G from G' and the
%         2-torsion point η = i(−1). A Théorème listing equivalent conditions
%         (i), (i bis), (ii) — existence of a polarising φ in Rep(G), then on
%         generators, then over R — and (iii), that the centraliser of C in
%         G'^0_R be a maximal compact subgroup. Closing: j satisfying a) and b)
%         is the same as j' : U → G'_R with j'(−1) = η.
%    6    G-subordinations: a ⊗-functor attaching to each V ∈ Rep(G) a Hodge
%         structure on V_C × S over an analytic space S, fibrewise polarisable.
%         Translated: a) a J'(s) : U → G'_R at each point satisfying the
%         theorem's conditions, b) real-analytic dependence on s.
%    7    the holomorphy condition needs a complex structure on the set J of
%         possible j; the functor S ↦ {G-subordinations} is then represented by
%         J, and Mor(S, J) is the set of rigidified subordinations. G(R) acts,
%         the orbits are open, and the connected components of J are
%         G(R)°/H(R) ∩ G(R)°.
%    8-9  the same without a base point, or rather with one: for S connected and
%         s ∈ S, π = π_1(S, ξ) acts, giving α : π → G(Q), and the universal
%         cover carries a G-subordination, giving f : S̃ → J; the pair (F, φ) is
%         equivalent to the pair (f, α) with f α-equivariant.
%   10-11 the base-point-free version: P_s = Isom_⊗(ω, Φ_s F) is a torseur under
%         G(Q); the P_s form a principal covering P of S with group Γ = G(Q),
%         and the category of G-subordinations is equivalent to that of pairs
%         (P, g) with g : P → J a Γ-morphism. Γ_0-rigidified subordinations for
%         Γ_0 commensurable with a G(Z); the moduli space is J/Γ_0, the
%         « variété de modules de niveau Γ_0 ». The run breaks off there.
%   13-14 the letter: the conjectures on Hodge structures he had told Mumford
%         contradict Tate's, which imply that a polarised analytic family over a
%         simply connected X is « algebraic » only if constant; and such
%         non-constant families are easy to produce over homogeneous spaces. Two
%         questions: can a single Hodge structure of even partial degrees be
%         algebraic, and is there a smooth projective surface Y with H^{1,1}
%         spanned by algebraic cycles and dim H^{0,2} ≥ 2?
%   15-16 the postscript: the data (G_R reductive connected over Q, i_0 : U_R →
%         G_R with i_0(−1) central and the centraliser K of i_0(√−1) maximal
%         compact), C_R = Centr(i_0), Q = G_R/C_R, and the « Siegel-space »
%         J = G(R)°/G(R)° ∩ C_R(R) ⊂ Q(R); J carries a complex structure
%         inherited from Q(R), every representation V with a lattice V_0 gives
%         an analytic family of Hodge structures on J, Γ acts, and the modular
%         varieties are the J/Γ with Γ commensurable with G(Z). A marginal note
%         records that J is not in general riemannian symmetric but fibres over
%         such a space with projective homogeneous fibres under K_C, hence is
%         simply connected.
%   17-18 the algebro-geometric reading: start from (V, φ) over Q with
%         G = SO(φ) and any i_0 with i_0(−1) = ±id; J generalises Siegel's
%         half-plane. Then the difficulty: taking G(R) compact, J/Γ = J is a
%         projective homogeneous space X under G_C by Borel, and the modular
%         family over it would, granting GAGA and the yoga of motives,
%         contradict Tate's conjectures, X being simply connected. « So I really
%         do not know what to believe! »
%   19    a necessary condition, granting Tate and Hodge, for J/Γ and the family
%         over it to be algebraic: for every g ∈ G(R)° and every Γ_0 of finite
%         index in G(Z), the smallest algebraic subgroup H of G with H_R ⊃
%         g i_0(U_R) g^{-1} and H ⊃ Γ_0 is G itself. A cancelled block; a
%         Question on Zariski density of G(Z); a question on Mumford's Boulder
%         talk.
%
% Notation fixed for the batch. Q, R, C, Z are set \mathbb, as he draws them
% with doubled strokes; so are his torus S (\mathbb{S}) and its norm-one
% subgroup U (\mathbb{U}), and the multiplicative group \mathbb{G}_{m}. The Tate
% objects are \mathbb{Q}(1), \mathbb{Q}(-1), \mathbb{Q}(-n) — his own notation.
% On pages 15-19 a roman Q is a different letter: the flag variety G_R/C_R, and
% Q(R) its real points; it is never the rationals there, which stay \mathbb{Q}.
% His script J, the set of the j's — the « Siegel-space » of the postscript — is
% \mathcal{J} throughout; \mathcal{J}/\Gamma is what he boxes. G' = \ker
% \varepsilon, and G^{0}, G'^{0} are the identity components as he writes them.
% \coprod is his own sign for the direct sum of the V^{p,q}, kept as drawn
% rather than normalised to \bigoplus. The two i's are distinguished as he does:
% i : \mathbb{G}_m → G in the French run, i_{0} : \mathbb{U}_{\mathbb{R}} →
% G_{\mathbb{R}} in the letter. The English of pages 13-19 stands as written and
% is not translated.
%
% What does not carry. The two halves of the batch are of very unequal
% legibility. The letter and its postscript (13-19) are written fair, at speed
% but without cancellation, and read almost end to end; the apparatus there is
% thin and concerns overwritten insertions and one block cancelled by a lattice
% of diagonal strokes (page 19). The French run (3-11) is the opposite: a fast
% palimpsest with insertions above the line, boxed afterthoughts, and long
% cancelled passages, in which the displayed formulas carry and the connective
% prose very often does not. Pages 6, 8 and 9 are the worst of it — page 6
% opens on a whole block struck through, and page 9's lower half is written over
% itself — and the apparatus is dense there in consequence.
%
% Batch boundary. The French run ends on page 11, on the « variété de modules de
% niveau Γ_0 », and the lower half of that leaf is blank: nothing crosses into
% batch 2 from it. The letter ends on page 19. Page 20 is a fresh cover for a
% new run, « Descriptions transcendantes de motifs en car. 0 … », which opens in
% batch 2; a \note says so. (Batch 2, having read them, records that leaves 21,
% 25, 32 and 35 are blank and that the rédaction itself begins on page 22.)

%
% Parallel pass, and what was reconciled afterwards. The folder's three batches
% were read concurrently, in three separate contexts, so batches 2 and 3 could
% not read the preceding batch's finished .tex to inherit its notation; each
% read the tail of the preceding facsimile instead. Two seams were reconciled
% in the orchestrating session once all three had landed, and nothing else was
% touched:
%   — batch 1 had set Q, R, C, Z, S, U, G_m in \mathbf and batches 2 and 3 in
%     \mathbb. Batch 1 was converted to \mathbb (205 occurrences), which is
%     both the corpus convention and the better record of the doubled strokes
%     he draws. No reading changed.
%   — the cover on page 20 had become a \section in batch 1 (closing it) and
%     again in batch 2 (opening it). It is now a heading in batch 2 alone,
%     where the run it names actually begins, and a \note in batch 1.
% Neither was a disagreement about what the leaves say.

\documentclass[11pt,a4paper]{article}
% Le préambule commun à toutes les transcriptions du fonds Grothendieck.
%
% Il tient en une idée : **l'incertitude se compose, elle ne se résout pas.**
% Une transcription qui lisse un mot douteux a détruit la seule chose pour
% laquelle on la faisait — un lecteur doit pouvoir savoir, à chaque endroit,
% ce qui était sur le feuillet et ce qui a été ajouté. Les six macros qui
% suivent sont donc l'appareil critique, et non de la décoration.
%
% Les mêmes commandes sont reconnues par `scripts/render.mjs`, qui produit la
% vue de lecture du site. Une macro ajoutée ici doit l'être aussi là-bas,
% faute de quoi le rendu échouera bruyamment — ce qui est voulu : un rendu
% silencieusement faux est pire qu'un rendu absent.

\ProvidesPackage{grothendieck}[2026/08/08 Transcriptions du fonds Grothendieck]

\usepackage{amsmath,amssymb,amsthm}
\usepackage{mathrsfs}
\usepackage{stmaryrd}
% Les diagrammes commutatifs sont partout dans ce fonds : c'est la langue
% même de la géométrie algébrique de Grothendieck, et une transcription qui
% les rendrait en prose ne transcrirait rien.
\usepackage{tikz-cd}
% Français par défaut — c'est la langue des feuillets — mais l'anglais est
% chargé aussi : la traduction et le résumé sont des documents anglais, et la
% typographie française (espaces avant les deux-points) y serait une faute.
% Ces documents basculent par \selectlanguage{english} en tête de corps.
\usepackage[english,french]{babel}
\usepackage{fontspec}
\usepackage{geometry}
\geometry{a4paper,margin=25mm,marginparwidth=28mm}
\usepackage{marginnote}
\usepackage{xcolor}
% ulem, not soul. `soul`'s \st and \ul are built on a character-by-character
% scan that XeLaTeX's font machinery breaks: the document fails with an
% undefined control sequence rather than degrading. ulem does the same two jobs
% and is Unicode-engine safe. `normalem` stops it hijacking \emph, which the
% transcriptions use for Grothendieck's own emphasis.
\usepackage[normalem]{ulem}
\usepackage{hyperref}
\hypersetup{colorlinks=true,linkcolor=black,urlcolor=black,pdfborder={0 0 0}}

% --- Métadonnées du lot -----------------------------------------------------
% Reprises telles quelles par le script de rendu, qui les affiche en tête de
% la vue de lecture. Elles fixent ce que le fichier prétend couvrir : sans
% elles, une transcription flotte sans point d'attache dans un fonds de seize
% mille feuillets.

\newcommand{\folder}[1]{\gdef\grothFolder{#1}}
\newcommand{\batch}[1]{\gdef\grothBatch{#1}}
\newcommand{\leaves}[2]{\gdef\grothFirst{#1}\gdef\grothLast{#2}}
\newcommand{\foldertitle}[1]{\gdef\grothFoldertitle{#1}}
\gdef\grothFolder{?}\gdef\grothBatch{?}\gdef\grothFirst{?}\gdef\grothLast{?}\gdef\grothFoldertitle{}

% --- Le résumé --------------------------------------------------------------
%
% Il ouvre la lecture modernisée, sous le titre « Résumé », et il est composé
% en petit. Ce n'est pas un ornement : c'est le seul passage du document qui
% ne soit pas des mathématiques, et un lecteur doit pouvoir le reconnaître
% avant de l'avoir lu — puis le sauter s'il connaît déjà le sujet, ou s'y
% arrêter s'il ne le connaît pas. Le filet qui le suit marque la reprise du
% corps.

\newenvironment{resume}
  {\small\setlength{\parskip}{0.45em}}
  {\par\medskip\hrule\medskip}

% --- Mots-clés ---------------------------------------------------------------
%
% En anglais, en clôture du résumé de la lecture modernisée : le vocabulaire
% moderne sous lequel chercher ce que les feuillets construisent. Le manifeste
% du site (`npm run manifest`) les extrait pour étiqueter la cote entière —
% cette ligne est la source unique des tags, il n'y a pas de fichier à côté.
\newcommand{\keywords}[1]{\par\smallskip\noindent
  {\footnotesize\textsc{Keywords} --- \textit{#1}}\par}

% --- L'appareil critique ----------------------------------------------------

% Le numéro de feuillet, en marge. C'est l'ancre qui permet de régler un
% désaccord sur une formule en regardant *un* feuillet plutôt qu'une chemise
% de six cents. Le site s'en sert aussi pour tourner les pages du fac-similé
% quand on fait défiler la transcription.
\newcommand{\leaf}[1]{%
  \marginnote{\footnotesize\color{gray!70}\textsf{#1}}%
  \label{leaf:#1}\ignorespaces}

% Illisible : on ne devine pas. Un mot inventé qui se lit comme les autres est
% le pire résultat possible.
\newcommand{\ill}{\textcolor{red!60!black}{[\,\ldots\,]}}

% Lecture douteuse : le mot est donné, et signalé comme incertain.
\newcommand{\uncertain}[1]{\uline{#1}}

% Ajout de l'éditeur — une lettre manquante, un mot sous-entendu.
\newcommand{\add}[1]{\textcolor{blue!50!black}{[#1]}}

% Biffé par Grothendieck lui-même. Ce qu'il a rayé fait partie du document :
% une rature dit souvent où la pensée a changé de direction.
\newcommand{\struck}[1]{\sout{#1}}

% Note du transcripteur, sur l'état matériel du feuillet.
\newcommand{\note}[1]{\par\smallskip{\small\color{gray!60!black}\textit{[#1]}}\par\smallskip}

% Note marginale de Grothendieck, distincte du corps du texte.
\newcommand{\marginal}[1]{\par\smallskip\noindent\hspace{1em}%
  {\small\color{gray!60!black}#1}\par\smallskip}

% --- Titre ------------------------------------------------------------------

\AtBeginDocument{%
  \begin{center}
    {\large\bfseries Fonds Alexandre Grothendieck --- cote n\textsuperscript{o}~\grothFolder}\\[2pt]
    {\itshape\grothFoldertitle}\\[4pt]
    {\small feuillets \grothFirst--\grothLast\ (lot \grothBatch)}\\[2pt]
    {\footnotesize\color{gray!60!black}Transcription établie d'après le fac-similé numérisé
     par l'Université de Montpellier.\\
     Les lectures incertaines sont soulignées, les illisibles notées [\,\ldots\,],
     les ajouts entre crochets.}
  \end{center}
  \vspace{1em}}

\endinput


\folder{13}
\batch{1}
\pages{1}{20}
\dating{1965}
\watermark{Édition de démonstration}
\foldertitle{Groupe de Galois motivique : notes manuscrites (s.d.), lettre (1965).}

\begin{document}

\section{Groupe de Hodge}

\note{Le titre est le sien : il l'inscrit, souligné, sur le feuillet 1, qui ne
porte rien d'autre que ces mots et, en dessous, « lettre à Mumford ». Ce
feuillet sert de chemise à tout ce qui suit et ne reçoit pas de numéro de page.
La suite court sans interruption du feuillet 3 au feuillet 11 ; il ne la
pagine pas.}

\page{3}

1) $G$ groupe algébrique sur $\mathbb{Q}$.

\[ \mathbb{S} = \prod_{\mathbb{C}/\mathbb{R}} \mathbb{G}_{m,\mathbb{C}} , \qquad \text{donc } \mathbb{S}(\mathbb{R}) \simeq \mathbb{C}^{*} . \]

Soit $\mathbb{U} \subset \mathbb{S}$ \add{$= \operatorname{Ker}(\text{Norme} : \mathbb{S} \to \mathbb{G}_{m,\mathbb{R}})$}

$\operatorname{Rep}(\mathbb{S}) \simeq$ catégorie des espaces vectoriels réels
$V_{\mathbb{R}}$, avec bigraduation sur
$V_{\mathbb{C}} = V_{\mathbb{R}} \otimes_{\mathbb{R}} \mathbb{C}$,

\[ V_{\mathbb{C}} \simeq \coprod_{p,q} V^{p,q} \]

satisfaisant les conditions

\[ V^{p,q} = \overline{V^{q,p}} . \]

\marginal{N.B. $\lambda \in \mathbb{C}^{*} \simeq \mathbb{S}(\mathbb{R})$ opère
sur $V^{p,q}$ par $\lambda^{p} \bar{\lambda}^{q}$}

\note{Le signe qu'il trace pour la somme directe des $V^{p,q}$ est un $\coprod$,
et il le garde d'un bout à l'autre du feuillet ; il est transcrit tel quel.}

L'homomorphisme canonique $\mathbb{G}_{m,\mathbb{R}} \xrightarrow{\ i\ }
\mathbb{S}$ correspond \uncertain{à} la graduation totale.

\[ \operatorname{Rep}(G) \xrightarrow{\ \text{oubli}\ } \operatorname{Modf}(\mathbb{Q}) \xrightarrow{\ \otimes_{\mathbb{Q}} \mathbb{R}\ } \operatorname{Modf}(\mathbb{R}) \xrightarrow{\ \otimes_{\mathbb{R}} \mathbb{C}\ } \operatorname{Modf}(\mathbb{C}) \]

La donnée sur le foncteur composé d'une bigraduation

\[ V_{\mathbb{C}} \simeq \coprod_{p,q} V_{\mathbb{C}}^{p,q} \]

satisfaisant $V^{p,q} = \overline{V^{q,p}}$

\marginal{i.e.\ d'une préstructure de Hodge, compatible avec la structure sur
$\mathbb{Q}$ donnée, et fonctorielle en l'objet de $\operatorname{Rep}(G)$}

équivaut à la donnée d'un hom.

\[ \mathbb{S} \xrightarrow{\ j\ } G_{\mathbb{R}} \]

Si on se donne un hom.\ \add{central}

\[ i : \mathbb{G}_{m} \longrightarrow G , \]

i.e.\ une graduation $V = \coprod_{n} V^{n}$ par le foncteur oubli $\omega$,
alors la condition

\[ V^{n}_{\mathbb{C}} = \coprod_{p+q=n} V^{p,q} \]

\page{4}

équivaut à la condition

\[ (\mathrm{a}) \qquad j \,|\, \mathbb{G}_{m,\mathbb{R}} = i_{\mathbb{R}} \]

Pour la suite, $i$ est supposé donné, \struck{$j$ \uncertain{compatible avec}
$i$ \ill{} \uncertain{au sens précédent}}, \add{$\mathbb{Q}(-1)$} objet
inversible (de $\operatorname{Rep}(G)$) à isom.\ près, i.e.\ un hom.

\[ \varepsilon : G \longrightarrow \mathbb{G}_{m} \]

On a alors, pour un $j$ donné, que $\mathbb{Q}(1)$ est de bidegré $(1,1)$ si et
seulement si on a les conditions

\[ (*) \qquad \varepsilon j(\lambda) = \lambda^{2} \ \text{ si } \lambda \in \mathbb{R}^{*} , \qquad \varepsilon(j(u)) = 1 \ \text{ si } u \in \mathbb{U} \]

i.e.

\[ (\mathrm{b}) \qquad \varepsilon j = \text{Norme} : \mathbb{S} \longrightarrow \mathbb{G}_{m,\mathbb{R}} \]

On supposera fixé $i$ et $\varepsilon$ satisfaisant \struck{\ill{}}
\add{$\varepsilon i(\lambda) = \lambda^{2}$}, et on s'intéresse aux $j$
satisfaisant a) et b).

Soit $j$ donné.

(Cherchons, pour $V \in \operatorname{Ob} \operatorname{Rep}(G)$ donné \add{de
degré $p$}, \ill{} structure de Hodge sur $V_{\mathbb{C}}$ (ou sur
$V_{\mathbb{R}}$), quelles sont les \struck{\ill{} de $R(G)$}
\struck{structures}

\[ V \times V \xrightarrow{\ \varphi\ } \mathbb{Q}(-p) \]

\struck{de Hodge, polarisées} qui définissent une \emph{polarisation} des
structures de Hodge correspondantes. Ce sont celles qui satisfont \struck{à la
condition que}

\[ (**) \qquad \varphi_{\mathbb{R}}(x, Cy) \]

forme \struck{symétrique} définie positive sur $V_{\mathbb{R}} \times V_{\mathbb{R}}$

\note{Le mot qui précède « structure de Hodge sur $V_{\mathbb{C}}$ » est d'un
seul jet et ne se laisse pas lire ; la parenthèse ouverte devant « Cherchons »
n'est jamais refermée sur le feuillet.}

\page{5}

où $C \in \mathbb{S}(\mathbb{R}) \simeq \mathbb{C}^{*}$ correspond à l'élément
$\sqrt{-1}$ de $\mathbb{C}^{*}$.

\marginal{[Ceci implique que $\varphi$ est \add{alternée si $n$ impair,}
symétrique si $n$ pair.]}

\marginal{N.B. Soit $G' = \operatorname{Ker} \varepsilon$, \struck{\ill{}} alors
$G \simeq G' \times_{\mathbb{Z}/2\mathbb{Z}} \mathbb{G}_{m,\mathbb{Q}}$, où
$\mathbb{Z}/2\mathbb{Z} \to G'$ est induit par $i$. Donc $G$ (avec $i$,
$\varepsilon$) se \uncertain{récupère} en connaissant $G'$ et son point d'ordre
$2$ $i(-1) = \eta$. La donnée de $i$ et $\varepsilon$ \ill{}}

\textbf{Théorème.} Pour $j$ donné, vérifiant les conditions a) et b), les
conditions suivantes \ill{} sont équivalentes et impliquent (iii) :

(i) Pour tout $V \in \operatorname{Ob} \operatorname{Rep}(G)$ \add{de degré
$n$}, [il existe une forme

\[ \varphi : V \otimes V \longrightarrow \mathbb{Q}(-n) \]

dans $\operatorname{Rep}(G)$ qui \struck{définit} une polarisation (pour la
structure de Hodge associée à $j$ sur $V$), i.e.\ on a :
\uncertain{$C = j(i)$}]

(i bis) Comme (i), mais \struck{seulement} pour des $V_{\alpha}$ donnés
\add{\ill{}} \struck{seulement}, \struck{\ill{}} \struck{que} l'hom.\
$G \to \operatorname{Aut}_{\mathbb{Q}}(V_{\alpha})$ soit injectif. De plus, on
suppose $G'$ réductif i.e.\ $G$ réductif.

(ii), \struck{(ii bis)} \struck{\ill{}} Comme (i), (i bis), mais avec $G$
remplacé par $G_{\mathbb{R}}$, et $\operatorname{Rep}(G)$ par
$\operatorname{Rep}(G_{\mathbb{R}})$. $G'_{\mathbb{R}}$ est réductif, et si
$K(j)$ est \struck{\ill{}} le centralisateur de $j(C)$ (intérieur --- $C$) dans
$G_{\mathbb{R}}$,

(iii) [Le \struck{\ill{}} centralisateur de l'image de $K(j)$ dans
$G^{0}_{\mathbb{R}} / i_{\mathbb{R}}(\mathbb{R}^{*+})$ est un sous-groupe
compact maximal, i.e.\ le centralisateur de $C$ dans ${G'}^{0}_{\mathbb{R}}$ est
un s-groupe compact \struck{\ill{}} maximal.]

\marginal{NB (iii) équivaut aux conditions \ill{} (iii bis) \ill{}
$G'' = G/\text{centre}$ \ill{} (iii ter) $\exists$ \ill{} $G$ \ill{}}

$j : \mathbb{S} \to G_{\mathbb{R}}$ \struck{\ill{}} satisfaisant a) et b)
équivaut à la donnée d'un hom.\ $j' : \mathbb{U} \to G'_{\mathbb{R}}$
satisfaisant $j'(-1) = \eta$

\note{La marge gauche de ce feuillet est écrite dans le sens de la hauteur et
n'est lisible que par fragments ; seules les étiquettes (iii bis) et (iii ter)
et le quotient $G/\text{centre}$ s'en dégagent.}

\page{6}

\struck{Pour tout espace analytique $S$, soit $F(S)$ l'ensemble \ill{} des
données suivantes \ill{} \struck{des} \struck{foncteurs} qui, à chaque
$V \in \operatorname{Rep}(G)$, associe \ill{}}

\note{Tout le haut du feuillet est un bloc encadré puis biffé en diagonale ;
il est repris immédiatement au-dessous, sous une forme qu'il garde.}

Soit $S$ un espace analytique \add{rigidifié} \struck{sur $S$}. Une
\emph{$G$-subordination} sur $S$ est la donnée d'un $\otimes$-foncteur qui, à
tout $V \in \operatorname{Ob} \operatorname{Rep}(G)$, associe une structure de
Hodge \add{mod.\ isogénie} sur

\[ V_{\mathbb{C}} \times S = V_{S} \]

[i.e.\ une bigraduation \struck{\ill{}} analytique réelle \struck{\ill{}} sur
$V_{S}$, $V_{S}^{p,q} = \overline{V_{S}^{q,p}}$, telle que la $p$-filtration
correspondante \ill{} filtration analytique complexe] \add{$V$ homogène de
degré $d$} et telle que $\forall\, s \in S$, la fibre en $s$ soit polarisable :
\struck{à l'aide d'une forme $\varphi : V \otimes V \to \mathbb{Q}(-n)$ \ill{}}
\struck{\ill{}}

\marginal{a) \ill{} de poids total \uncertain{cohérent} défini par $i$ ;
b) $\mathbb{Q}(-1)$ donne \ill{} de bidegré $(1,1)$, et}

\marginal{qui soit $G$-invariante (N.B. si $S$ est connexe, il suffit d'exiger
cette condition en un point de $S$, car \ill{})}

Traduisant \ill{} les conditions \ill{} la filtration, on trouve :

\begin{itemize}
  \item a) Pour $\forall\, V$ et $\forall\, s \in S$, on veut une façon
    \struck{de faire} \struck{fonctorielle} de mettre une structure de Hodge
    \add{(mod.\ isogénie)} polarisable sur les $V \in \operatorname{Ob}
    \operatorname{Rep}(G)$, i.e.\ donnée par un
    $j = J(s) : \mathbb{S} \to G_{\mathbb{R}}$ satisfaisant les conditions du
    th., i.e.\ $j' = J'(s) : \mathbb{U} \to G'_{\mathbb{R}}$ satisfaisant
    $j'(-1) = \eta$, et tel que
    $\operatorname{Zentr}_{{G'}^{0}_{\mathbb{R}}}(j'(\omega)(C))$ soit compact
    maximal.
  \item b) Pour $s$ variable, $J'(s)$ dépend de $s$ de façon analytique réelle.
\end{itemize}

\note{La prose qui relie les deux blocs encadrés de ce feuillet est écrite très
vite, par-dessus une première rédaction biffée, et ne se laisse pas suivre ;
seules les deux conditions a) et b) sont sûres.}

\page{7}

Enfin, pour exprimer la condition d'holomorphie de la fibration, il convient
d'introduire, \struck{sur} sur l'ens.\ $\mathcal{J}$ \add{des} $j$ (ou $j'$)
possibles, une structure de \emph{variété analytique} \struck{complexe}
naturelle, et la condition s'exprime par l'\emph{holomorphie} de l'application

\[ s \mapsto j(s) : S \longrightarrow \mathcal{J} . \]

\struck{En somme} Donc l'ens.\ des $G$-subordinations \struck{sur}
\add{rigidifiées} sur $S$ s'identifie à l'ens.\ $\operatorname{Mor}(S,
\mathcal{J})$ [morph.\ d'espaces analytiques]. Donc pour $S$ variable, on a un
foncteur représentable par $\mathcal{J}$.

N.B. $G(\mathbb{R})$ opère de façon évidente sur $\mathcal{J}$ (\ill{}) ; les
orbites sont ouvertes.

\struck{Soit $V_{0} \subset V$ un $\mathbb{Z}$-module de type fini engendrant
$V$. Alors $\forall$ \ill{} $G$-subordination sur $S$ définissant une structure
de Hodge (par un \ill{} $\otimes \mathbb{Q}$) sur $V_{0} \times S$, soit
$\Gamma \subset G(\mathbb{Q})$ le \ill{}}

Les composantes connexes de $\mathcal{J}$ sont de la forme

\[ G(\mathbb{R})^{\circ} / H(\mathbb{R}) \cap G(\mathbb{R})^{\circ} , \]

où $H$ est le centralisateur d'un $j : \mathbb{S} \to G$ dans $G$ ; ou encore

\[ G'(\mathbb{R})^{\circ} / H'(\mathbb{R}) \cap G'(\mathbb{R})^{\circ} , \]

\add{$H'$ le centralisateur d'un $j' : \mathbb{U} \to G'$ dans
$G'(\mathbb{R})^{\circ}$.}

\page{8}

Soit $C = \operatorname{Rep} G$ la $\otimes$-catégorie définie par $G$, et $S$
un espace analytique \emph{connexe}, $s \in S$. On veut classifier les
\struck{foncteurs} \add{$\otimes$-foncteurs} \add{(mod.\ isogénie)}

\[ F : C \longrightarrow \text{str. de Hodge} \]

\add{sur $S$}

[\struck{\ill{}} $\forall\, V \in \operatorname{Ob} C$, \add{$V$ de degré $n$,
$\exists\, \varphi : V \otimes V \to \mathbb{Q}(-n)$ telle que $F(\varphi)$ soit
une polarisation \add{fibre par fibre}} \add{[\ill{} une polarisation en un
point \ill{} de $S$, c'est Kif--kif].}]

\marginal{a) degré total (relatif à $i$) $=$ \ill{} degrés partiels ;
b) $\mathbb{Q}(-1)$ donne \ill{} str.\ de Hodge de bidegré $(1,1)$, et}

\marginal{\ill{} de Hodge \struck{\ill{}} s'appellent $G$-subordinations sur
$S$ (\ill{})}

\ill{} munis d'une structure fibrée

\[ C \longrightarrow \operatorname{Modf}(\mathbb{Q}) , \]

\ill{} on se donne un isomorphisme $\varphi$ entre cette donnée et le foncteur
oubli.

\struck{Alors} Soit $\pi = \pi_{1}(S, \xi)$. Alors $\pi$ opère \add{à g.} sur le
foncteur \struck{fbt} « $\mathbb{Q}$-fibrés sur $S$ » sur la catégorie des str.\
de Hodge sur $S$, d'où sur le foncteur composé avec $F$, d'où une \ill{} de
$(F, \varphi)$ :

\[ \alpha : \pi = \pi_{1}(S, \xi) \longrightarrow G(\mathbb{Q}) \]

Soit d'autre part $\widetilde{S}$ le revêtement universel de $S$ pointé
au-dessus de $s$ par $\widetilde{s}$. Alors la donnée de $\varphi$ permet de
considérer le foncteur

\[ V \longmapsto F(V) \longmapsto F(V)|_{\widetilde{S}} \]

comme une $G$-subordination sur $\widetilde{S}$. Donc on trouve un \ill{}

\note{Ce feuillet reprend, sans point-base d'abord puis avec, la construction du
feuillet 7 ; les deux encadrés de la marge gauche reprennent les conditions a)
et b) du feuillet 6.}

\page{9}

\struck{deuxième} élément de structure

\[ f : \widetilde{S} \longrightarrow \mathcal{J} \]

Les deux éléments de structure sont reliés par la condition que $f$ est
compatible avec les actions de $\pi$, quand on le fait opérer sur $\mathcal{J}$
via $\alpha : \pi \to G(\mathbb{Q}) \subset G(\mathbb{R})$ et l'action
\add{à droite} de $G(\mathbb{R})$ sur $\mathcal{J}$ (\ill{} à l'intérieur
gauche).

La donnée d'un couple $(F, \varphi)$ équivaut à celle d'un couple
$(f, \alpha)$ satisfaisant les deux conditions précédentes.

\marginal{Il faudrait fixer des \emph{réseaux}, et des sous-groupes discrets
commensurables de $G(\mathbb{Q})$, du type $G(\mathbb{Z})$ \ill{}}

Autre façon de présenter les choses, indépendante du choix d'un point-base de
$S$. \struck{\ill{}}

Notons que les foncteurs fibres \add{$\Phi_{s}$} sur \struck{str.\ de Hodge}
$\operatorname{Hodge}(S)$ définis par les divers points $s \in S$ \ill{}
\add{$\Phi_{s} \circ F$}. On dit \ill{} que \ill{} sur $C =
\operatorname{Rep}(G)$ si les $\Phi_{s} \circ F$ sont isomorphes
\struck{\ill{}} \ill{}. Soit donc, pour

\note{La moitié basse de ce feuillet est écrite par-dessus une première
rédaction et ne se laisse pas relever ; la phrase se poursuit sur le feuillet
10, où elle redevient lisible.}

\page{10}

$s \in S$, $P_{s}$ le torseur sous $G(\mathbb{Q})$

\[ P_{s} = \operatorname{Isom}_{\otimes}(\omega, \Phi_{s} F) \]

Pour $s$ variable, les $P_{s}$ forment \add{(à droite)} \ill{} un revêtement
principal $P$ de $S$, de gr.\ $\Gamma = G(\mathbb{Q})$. \uncertain{On} \ill{}
fait ce qu'il fallait pour que le composé

\[ C = \operatorname{Rep}(G) \xrightarrow{\ F\ } \operatorname{Hodge}(S) \xrightarrow{\ V \mapsto V \times_{S} P\ } \operatorname{Hodge}(P) \]

\struck{\ill{}} définisse sur $P$ une $G$-subordination, d'où un morphisme

\[ g : P \longrightarrow \mathcal{J} \]

\struck{Ceci posé,} Ce dernier est \add{compatible} (opérant sur $\mathcal{J}$
à \emph{droite}) avec l'action de $\Gamma = G(\mathbb{Q})$. [Et la catégorie des
\struck{$G$-foncteurs} \struck{$F$} $G$-subordinations $F$ sur $S$ compatible
avec le foncteur \ill{} est équivalente à celle des couples $(P, g)$, $P$ un
rev.\ principal de $S$ de gr.\ $\Gamma$, $g : P \to \mathcal{J}$ un
$\Gamma$-morphisme.]

Soit $\Gamma_{0} \subset \Gamma = G(\mathbb{Q})$ un s-groupe [\ill{},
commensurable à un $G(\mathbb{Z})$]. On appelle \emph{$G$-subordination
$\Gamma_{0}$-rigidifiée} sur $S$ la donnée d'un rev.\ principal $P_{0}$ de $S$
de

\page{11}

groupe $\Gamma_{0}$, et d'un $\Gamma_{0}$-morphisme de $P_{0} \to \mathcal{J}$.

Si $\Gamma_{0} = \{e\}$, on trouve la notion de $G$-subordination rigidifiée ;
si $\Gamma_{0} = G(\mathbb{Q})$, la notion de $G$-subordination
\struck{\ill{}} tout court. On peut expliciter la notion par la donnée d'un
$F$ \add{compatible avec $\omega$,} et d'une \struck{\ill{}}
$\Gamma_{0}$-rigidification de $F$, i.e.\ une section de $P_{0}$ \ill{} pour la
rigidification naturelle de $\Gamma_{0}$ \ill{}

Pour cette structure, on a un espace modulaire, qui est \ill{}
($\mathcal{J}$, avec action de $\Gamma_{0}$ à droite). \ill{} la question

\[ \mathcal{J} / \Gamma_{0} \]

[avoir une structure d'espace analytique quand $\Gamma_{0}$ agit proprement
\ill{} et discret] et pour cela d'une \emph{variété de modules}
\struck{\ill{}} de niveau $\Gamma_{0}$.

\note{Le feuillet s'arrête là, et sa moitié basse est blanche : la suite de
l'argument ne passe pas dans la deuxième liasse.}

\section{Lettre à Mumford}

\note{Ces mots sont les siens, inscrits sous « Groupe de Hodge » sur le feuillet
1 qui sert de chemise. La lettre occupe les feuillets 13 et 14 ; son
post-scriptum, qui porte l'essentiel des mathématiques, court des feuillets 15
à 19. Elle est en anglais et le reste ici.}

\page{13}

Algiers Nov.\ 1, 1965

Dear Mumford,

I have become aware that I have been a bit rash with the conjectures I told you
about Hodge structures, as in the form I stated them they contradict Tate's
conjecture, which implies the following: if $X$ is a smooth, connected simply
connected scheme over $\mathbb{C}$, and $V$ a polarized complex analytic family
of Hodge structures parametrized by $X^{\mathrm{hol}}$, then $V$ is
« algebraic » only if it is a constant family. Now it is easy to get examples of
$V$'s, with $X$ say \struck{any} homogeneous space under an affine group over
$\mathbb{C}$)

\page{14}

which are non constant, therefore (one hopes!) not « algebraic » (as one would
like Tate's conjectures to hold).

However, one should try to test if one really cannot get an algebraic $V$ that
way. To start, I wonder if a single Hodge-structure $V$ with even partial
degrees (therefore giving rise to a « Hodge-\ill{} » $G$ s.th.\ $G_{\mathbb{R}}$
is compact) can be algebraic. \add{(except when $G$ is commut.)}

For instance, do you know an algebraic smooth projective surface $Y$ over
$\mathbb{C}$, such that $H^{1,1}(Y, \mathbb{C})$ is spanned by algebraic cycles,
and $\dim H^{0,2}(Y, \mathbb{C}) \geqslant 2$ ?? I would appreciate your
comments!

Yours

A Grothendieck

P.S. What about your $\Theta$-fns, are you going to write a book?

\note{Le mot composé « Hodge-… » est recouvert par la parenthèse insérée
au-dessus de la ligne et n'est pas récupérable ; la parenthèse « $G$ s.th.\
$G_{\mathbb{R}}$ is compact » est bien celle du feuillet, écrite un cran plus
bas que le reste de la phrase.}

\page{15}

P.S. Here what are the modular varieties one gets from consideration of
motive-theoretic galois theory.

$G$ $=$ reductive connected alg.\ group defined over $\mathbb{Q}$.

\note{L'indice $\mathbb{R}$ de ce premier $G$ est raturé ; il écrit ensuite
$G_{\mathbb{R}}$ partout.}

$\mathbb{U}_{\mathbb{R}} =$ dimension one twisted real Torus, with
$\mathbb{U}_{\mathbb{R}}(\mathbb{R})$ identified with the group of complex nbs
of module $1$

$i_{0} : \mathbb{U}_{\mathbb{R}} \longrightarrow G_{\mathbb{R}}$ homomorphism of
algebraic groups over $\mathbb{R}$, s.th.\ $i_{0}(-1)$ is central, and s.th.\
the centralizer $K_{\mathbb{R}}$ of $i_{0}(\sqrt{-1})$ is such that
$K_{\mathbb{R}}(\mathbb{R})^{\circ}$ is maximal compact in
$G_{\mathbb{R}}(\mathbb{R})^{\circ}$. Moreover, for every $\mathbb{R}$-simple
component of $G_{\mathbb{R}}$, the corresponding component of $i_{0}$ is not
trivial.

\[ C_{\mathbb{R}} = \operatorname{Centr}_{G_{\mathbb{R}}}(i_{0}) \subset K_{\mathbb{R}} \]
\[ Q = G_{\mathbb{R}} / C_{\mathbb{R}} \qquad \text{(a connected component)} \]

« Siegel-space » $=$

\[ \mathcal{J} = G(\mathbb{R})^{\circ} / G(\mathbb{R})^{\circ} \cap C_{\mathbb{R}}(\mathbb{R}) \subset Q(\mathbb{R}) \]

\note{Le $Q$ de $Q = G_{\mathbb{R}}/C_{\mathbb{R}}$ et de $Q(\mathbb{R})$ n'est
pas le corps des rationnels : c'est la variété de drapeaux. La distinction est
tenue partout dans ce post-scriptum.}

N.B. $\mathcal{J}$ has a natural complex structure (inherited from the one on
$Q(\mathbb{R})$ in fact) invariant by operation of $G(\mathbb{R})$. Every linear
representation (over $\mathbb{Q}$) of $G$ into a vector space $V$ (over
$\mathbb{Q}$), together with a lattice $V_{0}$ of $V$, defines a « complex
analytic family of Hodge structures » on $\mathcal{J}$, at least when one fixes
a « total degree » $n$

\marginal{(\struck{\ill{}} \add{necessarily} globally polarizable --- to get a
polarization of the family, one picks a « form » of polarization on $V$
invariant by $G$).}

\page{16}

and assumes that $i_{0}(-1)$ operates on $V$ as \add{homothety} $(-1)^{n}$. (It
is not hard to describe $\mathcal{J}$ by universal properties, as
\struck{\ill{}} representing some functor $F(X)$, $X$ \struck{complex analytic},
expressing that every $V$ \add{as above} should \struck{\ill{}} define, in an
additive and multiplicative way, for varying $V$, a Hodge-structure $V$ over
$X$).

\marginal{Note that in general $\mathcal{J}$ is \emph{not} riemannian symm., but
is a fiber space over a R.S. space, with fibers which are projective complex
homogeneous spaces under $K_{\mathbb{C}}$. This implies at least that
$\mathcal{J}$ is \emph{simply connected}, hence is the univ.\ covering of
$M = \mathcal{J}/\Gamma$ if $\Gamma$ operates freely.}

If $\Gamma \subset G(\mathbb{Q})$ is the group leaving $V_{0}$ fixed, then
$\Gamma$ operates on $(V, \mathcal{J})$. $\mathcal{J}$ expects to carry
algebraic (and usually, $\mathbb{Z}$) structures.

The modular varieties one gets from consideration of Hodge-structures
\emph{over} $\mathbb{Q}$ are the varieties

\[ \mathcal{J} / \Gamma , \]

where $\Gamma \subset G(\mathbb{Q}) \cap G(\mathbb{R}^{\circ})$ is
commensurable with $G(\mathbb{Z})$ [for some \struck{\ill{}} matrix
representation of $G$, \struck{\ill{}} giving a meaning to $G(\mathbb{Z})$].
However, it is possible that one will have to impose on the data some extra
conditions, either arithmetic (involving existence of « Frobenius elements » in
\struck{\ill{}} $I(\mathbb{Q})$, $I = G/\text{int.\ aut.}$), or a geometric
\add{(on $G_{\mathbb{R}}$)} (much \ill{} : $G_{\mathbb{R}}$ has \emph{no compact
factor}). A typical case which fits in the general \struck{picture}
\add{description}, but for which I have no evidence so far if $\mathcal{J}/
\Gamma$ has

\page{17}

an algebro-geometric interpretation in terms of \add{modular} (motives, is the
following:

start with a vector space of finite dimension $V$ over $\mathbb{Q}$, with
\struck{\ill{}} a fundamental bilinear form $\varphi$ given, symmetric or
alternating, take $G = SO(\varphi)$, take any

\[ i_{0} : \mathbb{U}_{\mathbb{R}} \longrightarrow G_{\mathbb{R}} \]

s.th.\ $i_{0}(-1) = \mathrm{id}$ (if $\varphi$ symmetric), $i_{0}(-1) =
-\mathrm{id}$ (if $\varphi$ skew-symmetric), \struck{\ill{}} \add{and a
positivity} condition of $\varphi$ relative to the bigrading of $V_{\mathbb{C}}$
defined by $i_{0}$, which \struck{I will} not write down.

Then $\mathcal{J}$ is a direct generalization of Siegel's half-plane, and
classifies the Hodge-structures \struck{\ill{}} with underlying
$\mathbb{Q}$-vector space $V$, \add{i.e.} \add{M} bigradings of $V_{\mathbb{C}}$,
\struck{\ill{}} \add{satisfying} the usual symmetry condition for complex
conjugation, which are compatible with $\varphi$ (i.e.\ $\varphi$ is a
\add{polarized} polarization), and for which the $\dim V_{\mathbb{C}}^{p,q}$
have given values (depending on the weights of \ill{} question on
$V_{\mathbb{C}}$ \ill{}). The thought is that, \struck{\ill{}} of these special
types of $\mathcal{J}$, $\mathcal{J}/\Gamma$ are in a way « universal » for all
others

\page{18}

(namely the $V$'s considered above are induced from one's on such special types
of $\mathcal{J}$), if the « modular family » $V$ \add{such a} of Hodge
structures on ($\mathcal{J}/\Gamma$ was) « algebraic » i.e.\ came from a
relative complex analytic space over the corresponding \ill{} $\mathcal{J}/
\Gamma$, \add{(or even on some underlying structure of alg.\ var.\ over
$\mathbb{C}$)}, [the same would be true for \emph{every} modular family on some
$\mathcal{J}/\Gamma$.]

Taking the \uncertain{case} \struck{$G(\mathbb{R})$} $G(\mathbb{R})$ compact,
\struck{this} and thus

\[ \mathcal{J}/\Gamma = \mathcal{J} = \text{complex projective homogeneous space } X \text{ under } G_{\mathbb{C}} \text{ by Borel,} \]

we would get a \add{polarized} motive over the analytic space $\mathcal{J}$
corresponding to $X$. If we admit the yoga \add{GAGA,} (coming from
\add{polarized} (abelian varieties)) that such a family must also be a motive
over $X$, \add{(or the above strengthening of the original assumption)}, [we
\struck{would} get a contradiction with] Tate's conjectures $\leftarrow$
(because $X$ is simply connected, as well known). So I really do not know what
to believe!

\page{19}

To come back to the general case, maybe I should add that a \emph{necessary}
condition for $\mathcal{J}/\Gamma$ to be algebraic and the Hodge-structure
$V/\Gamma$ over $\mathcal{J}/\Gamma$ algebraic too, is (\emph{granting} Tate's
and Hodge's conjectures): for every $g \in G(\mathbb{R})^{\circ}$, and every
$\Gamma_{0}$ of finite index in $G(\mathbb{Z})$, the smallest \add{alg.}
subgroup $H$ of $G$ such that $H_{\mathbb{R}}$ contains
$g\, i_{0}(\mathbb{U}_{\mathbb{R}})\, g^{-1}$ \emph{and} $\Gamma_{0}$ is $G$
itself.

\struck{I did not try to test directly this condition, even in the case of the
usual Siegel modular space; if it should be false in that case, \ill{} we know
that $\mathcal{J}/\Gamma$ is algebraic and that so are the $V/\Gamma$ over it,
it would follow that either Hodge's or Tate's conjecture is wrong.}

\note{Ce bloc est annulé par un quadrillage de traits diagonaux ; il se lit
néanmoins presque entièrement et est transcrit. Le « not » de « I did not try »
est lui-même biffé à l'intérieur du bloc, d'un trait plus épais.}

Question: If $G_{\mathbb{R}}$ is semi-simple without compact factors, is
$G(\mathbb{Z})$ always Zariski-dense in $G$ ??

By the way, did you make out if for the modular varieties $\mathcal{J}/\Gamma$
of your Boulder talk (which I just read), the image of $\mathcal{J}/\Gamma$ in
the usual Siegel modular variety is an algebraic variety, or at least
constructible?

\note{Le feuillet 20 ne porte rien d'autre que ces mots de sa main :
« Descriptions transcendantes de motifs en car.\ 0 et de catégories de motifs
en car.\ 0 » et, sous un trait, « doubles réseaux etc ». C'est la chemise de
la rédaction qui suit, laquelle ne commence pas dans cette liasse : elle ouvre
la deuxième, où ce titre est repris en tête de section.}

\end{document}
